\section*{Hoofdstuk \ref{ch:b2:corrosion} --- Corrosie en bescherming}

\begin{solution}{exo:b2:corrosion:1}
$j = \qty{0.020}{A/m^2}$: $de/dt = 0.020 \times 65.38/(2 \times 96\,485 \times 7.13
\times 10^6) = \qty{9.5e-13}{m/s}$, dus \qty{0.030}{mm} per jaar.
\end{solution}

\begin{solution}{exo:b2:corrosion:2}
Stalen bouten op koper: ijzer, het minst edele metaal, is de anode, en de kleine bouten
op een grote koperen kathode corroderen snel. Zinken nagels in staal: het zink corrodeert
en beschermt het staal eromheen.
\end{solution}

\begin{solution}{exo:b2:corrosion:3}
In de smalle spleet tussen de platen en onder de boutkop, waar het water
traag ververst wordt en zijn zuurstof snel opgebruikt is: door \omterm{def:b2:corrosion:differential}{differentiële beluchting}
worden deze ingesloten zones anoden, en de open oppervlakken kathoden.
\end{solution}

\begin{solution}{exo:b2:corrosion:4}
Het conservenblik: tin, edeler dan ijzer, maakt van het blootgelegde staal de anode van een
grote kathode. Op de gegalvaniseerde emmer corrodeert het zink in plaats van het
staal.
\end{solution}

\begin{solution}{exo:b2:corrosion:5}
Roeren maakt de \omterm{def:b2:current-potential-curves:limiting-current}{diffusielaag} dunner: het zuurstofplateau stijgt, de anodische kromme
van ijzer snijdt het hoger, zodat zowel de \omterm{def:b2:corrosion:uniform}{corrosiepotentiaal} als de \omterm{def:b2:corrosion:uniform}{corrosiestroom}
toeneemt. Zonder zuurstof blijft in neutraal water als enige kathodische reactie de
trage reductie van water over: de \omterm{def:b2:corrosion:uniform}{corrosiepotentiaal} daalt en de
\omterm{def:b2:corrosion:uniform}{corrosiestroom} wordt zeer klein.
\end{solution}

\begin{solution}{exo:b2:corrosion:6}
$t = 0.90 \times 5000 \times 2 \times 96\,485/(65.38 \times 0.10) = \qty{1.33e8}{s}$,
ongeveer 4.2~jaar.
\end{solution}

\begin{solution}{exo:b2:corrosion:7}
Lading per kilogram, $nF/M$: zink \qty{820}{A.h/kg}, magnesium \qty{2.21e3}{A.h/kg},
aluminium \qty{2.98e3}{A.h/kg}. Standaardpotentialen: \ce{Zn^2+}/\ce{Zn} $-0.76$,
\ce{Mg^2+}/\ce{Mg} $-2.36$, \ce{Al^3+}/\ce{Al} \qty{-1.68}{V}. In een droge bodem met hoge weerstand
moet de stroom door een grote weerstand gedreven worden: magnesium, dat veel verder
onder ijzer ligt, biedt de grootste drijvende spanning.
\end{solution}

\begin{solution}{exo:b2:corrosion:8}
$U = 2.0 \times 4.0 + 1.5 = \qty{9.5}{V}$; vermogen \qty{19}{W}; per jaar $19 \times
8766 = \qty{1.7e5}{W.h}$, \qty{1.7e2}{kW.h}.
\end{solution}

\begin{solution}{exo:b2:corrosion:9}
Op de waterlijn wordt het metaal bevochtigd door zuurstofrijk water: een kathode; net
eronder is het water minder belucht en wordt de zone de anode van de cel die
met de waterlijn gevormd wordt: daar corrodeert het staal, in een band.
\end{solution}

\begin{solution}{exo:b2:corrosion:10}
In belucht water valt de \omterm{def:b2:corrosion:uniform}{corrosiepotentiaal} van roestvast staal op zijn
passieve plateau: een minuscule stroom, en de film herstelt zichzelf. Chloride-ionen breken
de film op zwakke plekken; het blote metaal is daar een kleine anode tegenover een
enorme passieve kathode, zodat zijn stroomdichtheid enorm is. In de put
hydrolyseren de metaalionen en verzuren ze de oplossing, wordt chloride aangetrokken om
de lading te compenseren, en kan de film zich niet opnieuw vormen: de put wordt dieper.
\end{solution}

\begin{solution}{exo:b2:corrosion:11}
Zuurstof wordt gereduceerd op alle \qty{22}{cm^2}: \omterm{def:b2:current-potential-curves:ie-curve}{kathodische stroom} $22 \times 5 =
\qty{110}{\micro A}$, volledig geleverd door de oxidatie van de \qty{2}{cm^2} staal
bij de verbinding, $j = \qty{55}{\micro A/cm^2}$. Het verlies is $5.5$ maal dat van ijzer
bij \qty{10}{\micro A/cm^2}, ongeveer \qty{0.64}{mm} per jaar.
\end{solution}

\begin{solution}{exo:b2:corrosion:12}
IJzer is immuun onder de lijn \ce{Fe^2+}/\ce{Fe}, $-0.41 + 0.0296\log 10^{-6} =
\qty{-0.59}{V}$: onder ongeveer \qty{-0.6}{V} gehouden lost het staal niet meer op. Veel
lager wordt water in toenemende mate op het staal gereduceerd (de lijn \ce{H+}/\ce{H2} ligt bij pH 7 op
\qty{-0.41}{V} en de \omterm{def:b2:current-potential-curves:overpotential}{overspanning} op staal bedraagt enkele tienden volt):
er gaat stroom verloren en waterstof doet de deklaag loslaten.
\end{solution}

\begin{solution}{pb:b2:corrosion:1}
\textbf{1.} \ce{Fe -> Fe^2+ + 2e-}; \ce{O2 + 2H2O + 4e- -> 4OH-}.
\textbf{2.} De reductie van zuurstof, traag en op haar diffusieplateau,
begrenst de \omterm{def:b2:current-potential-curves:ie-curve}{kathodische stroom}; de anodische kromme van ijzer snijdt haar op dat plateau.
\textbf{3.} $j = \qty{0.10}{A/m^2}$: $0.10 \times 3.156 \times 10^7/(2 \times 96\,485)
\times 55.845 = \qty{913}{g}$ per vierkante meter en per jaar.
\textbf{4.} $913/7.87 \times 10^6 = \qty{1.16e-4}{m}$, \qty{0.12}{mm} per jaar.
\textbf{5.} $4/0.116 = 34$~jaar.
\textbf{6.} De corrosie is niet uniform: spleten, verschillen in beluchting
tussen bodems, en gebreken in de deklaag concentreren de \omterm{def:b2:current-potential-curves:ie-curve}{anodische stroom} op
kleine oppervlakken, waar putten de wand lang voordien doorboren.
\textbf{7.} $E^\circ = \Delta_f G^\circ/(nF)$ voor elk koppel: \ce{Fe} $-0.41$, \ce{Zn} $-0.76$,
\ce{Mg} $-2.36$, \ce{Al} \qty{-1.68}{V}.
\textbf{8.} Zink, magnesium en aluminium, die allemaal onder ijzer liggen.
\textbf{9.} Aluminium passiveert: zijn oxidefilm stopt het oplossen en het
levert dan weinig stroom, tenzij het gelegeerd is om de film te verhinderen (wat werkt
in zeewater, niet in bodems).
\textbf{10.} De stroom moet de weerstand van de bodem overwinnen: de drijvende
spanning tussen de anode en het beschermde staal bedraagt ongeveer \qty{1.9}{V} voor
magnesium tegenover \qty{0.35}{V} voor zink (uit de standaardpotentialen).
\textbf{11.} $3.156 \times 10^7 \times 24.305/(2 \times 96\,485 \times 0.50) =
\qty{7.95}{kg}$ per ampère en per jaar.
\textbf{12.} $\pi \times 0.30 \times 10\,000 = \qty{9.42e3}{m^2}$.
\textbf{13.} \qty{94.2}{m^2}.
\textbf{14.} $0.020 \times 94.2 = \qty{1.88}{A}$.
\textbf{15.} $1.885 \times 20 \times 3.156 \times 10^7 = \qty{1.19e9}{C}$.
\textbf{16.} $1.19 \times 10^9 \times 24.305/(2 \times 96\,485 \times 0.50) =
\qty{3.00e5}{g}$, \qty{300}{kg} (of $1.885 \times 20 \times 7.95$).
\textbf{17.} $300/15 = 20$ anoden.
\textbf{18.} De weerstand van de bodem beperkt hoe ver één anode haar stroom
langs de buis kan duwen: verspreide anoden beschermen de buis gelijkmatig.
\textbf{19.} Een gelijkrichter stuurt stroom van een inerte anode in een anodebed,
door de bodem, naar de buis, die met zijn negatieve pool verbonden is.
\textbf{20.} $1.885 \times 5.0 + 1.5 = \qty{10.9}{V}$.
\textbf{21.} $10.9 \times 1.885 = \qty{20.6}{W}$; $20.6 \times 8766 = \qty{1.8e5}{W.h}$,
\qty{180}{kW.h} per jaar.
\textbf{22.} Water wordt aan de buis gereduceerd: er ontstaat waterstof, die de deklaag loslicht en
het staal bros kan maken, en er gaat stroom verloren.
\textbf{23.} $E < -0.41 + 0.0296\log(10^{-6}) = \qty{-0.59}{V}$.
\textbf{24.} Twintig jaar bescherming vereist $\boldsymbol{\approx \qty{3.0e2}{kg}}$
magnesiumanoden.
\end{solution}
